%% fig_heatmap.tex
%% 分離の「動作マップ」: パラメータ平面 (|Delta alpha|E/p, beta pE) 上での
%% 平均速度 v=L/T1 を上から見たヒートマップで表す。同じ動作点でも
%% m（棒長）で v が大きく異なれば、速度（移動度）の差として
%% 電気的・幾何的性質を分離できることを示す。
\begin{tikzpicture}
  \begin{groupplot}[
      group style={group size=2 by 1, horizontal sep=4.1cm},
      width=4.25cm, height=4.35cm, scale only axis,
      view={0}{90},                      % 真上から見た 2D ヒートマップ
      tick align=outside,
      tick label style={font=\footnotesize}, label style={font=\small},
      title style={font=\small},
      xlabel={$\frac{|\Delta \alpha|E}{p}$},
      xtick={-1.585,0,1,1.585,2.322,3.322},
      xticklabels={$\tfrac13$,$1$,$2$,$3$,$5$,$10$},
      ylabel={$\beta pE$}, ytick={0.5,1,2,3},
      xmin=-1.585, xmax=3.322, ymin=0.5, ymax=3,
      enlargelimits=false,
      colormap/viridis,
      colorbar, colorbar style={font=\footnotesize, width=0.22cm,
        ylabel={$v=L/T_1$}},
    ]
    \nextgroupplot[title={$m=3$}, point meta min=0.25, point meta max=1.0]
      \addplot3[surf, shader=interp, mesh/cols=6]
        table[x=log2delta, y=beta, z=v]{\figdata/heat_m3.dat};
    \nextgroupplot[title={$m=6$}, point meta min=0.13, point meta max=0.62]
      \addplot3[surf, shader=interp, mesh/cols=6]
        table[x=log2delta, y=beta, z=v]{\figdata/heat_m6.dat};
  \end{groupplot}
\end{tikzpicture}
